\documentclass[11pt,letterpaper]{article}
\usepackage[T1]{fontenc}
\usepackage[utf8]{inputenc}
\usepackage[scaled=0.95]{helvet}
\renewcommand{\familydefault}{\sfdefault}
\usepackage[top=1.8cm,bottom=1.6cm,left=2.2cm,right=2.2cm]{geometry}
\usepackage{microtype}
\usepackage[hidelinks]{hyperref}
\pagestyle{empty}
\setlength{\parindent}{0pt}
\setlength{\parskip}{8pt}

\begin{document}

{\large\bfseries Grover Aruquipa}\par\vspace{-6pt}
{\small Qu\'ebec, QC, Canada \ \textbullet\ \ grover.aruquipa.9@gmail.com \ \textbullet\ \ groveraruquipa.github.io}

\vspace{6pt}
Hiring Team, Robotics Center of Excellence, Rockwell Automation\\
Robotics Researcher (R26-4200), Waterloo

Dear Hiring Team,

My first robot was a Delta for picking products, built for my bachelor's thesis in La Paz. I then spent two years in industry on the data side of automation, doing machine learning on circuit-breaker sensor data at Hitachi ABB Power Grids R\&D in Zurich and computer vision for anomaly detection on production lines at Sencon. I came back to research because I wanted to understand the robot itself, and now that my PhD is ending I want to bring both halves back to the factory floor. A team whose stated goal is to see code moving steel in a real plant is exactly where I want to do that.

For the last three years at Universit\'e Laval, with Prof.\ Cl\'ement Gosselin, I have worked on robotic grasping and manipulation. I designed, built and controlled three manipulators with 4, 7 and 8 actuators in which the gripper is driven through the mechanism, so the robot can rotate a grasped object without limit. I wrote the 500\,Hz real-time control in C++ and validated a coaxial wrist-gripper over ten continuous revolutions with sub-degree tracking while it held the part. I also wrote a review of robotic grasping published in the ASME Journal of Mechanisms and Robotics, and at JKU Linz I developed grasp-quality indices with Prof.\ Andreas M\"uller. Right now I am building a Delta robot with an integrated gripper that estimates grasp force from motor currents, to handle fragile objects at speed.

I want to be honest about the gap. I have not trained VLA models or run large-scale reinforcement learning, and my simulation work so far has been in MuJoCo and Simscape rather than Isaac Lab. What I bring is the other end of the pipeline: I know why a grasp fails, what the contact forces and the kinematics allow, and how to make a policy's output survive on real motors. My PyTorch and TensorFlow experience comes from industry, and closing the gap on foundation models and sim-to-real is the part of this job I am most eager to do.

I am based in Qu\'ebec, I defend my PhD in January 2027, and I could start in Waterloo after that, on site on the days your hybrid policy asks for. After three years here I would like to keep building my career in Canada, and Rockwell is the company I would most like to do it with. Prof.\ Gosselin, Prof.\ Patrice Lambert and Prof.\ M\"uller can speak to my work; their contact details are in my CV.

Thank you for your time. I would be glad to show you the robots running.

Sincerely,\\
Grover Aruquipa

\end{document}
